\documentclass[12pt]{article}
% Adaptado de C:/GitHub/Mestrado/Template/Template.tex.
\usepackage[T1]{fontenc}
\usepackage[utf8]{inputenc}
\usepackage[brazilian]{babel}
\usepackage{graphicx,booktabs,array,amsmath,amssymb,microtype}
\usepackage{geometry}
\usepackage{pdflscape,caption}
\geometry{a4paper,margin=2.5cm}
\usepackage{parskip,csquotes}
\interfootnotelinepenalty=10000
\usepackage[backend=biber,style=abnt]{biblatex}
\addbibresource{references.bib}
\usepackage{hyperref}
\hypersetup{colorlinks=true,urlcolor=blue,linkcolor=blue,citecolor=blue}
\newcommand{\instituicao}{UNIVERSIDADE FEDERAL DE SANTA CATARINA}
\newcommand{\departamento}{Programa de Pós-Graduação em Economia}
\newcommand{\curso}{Mestrado}
\newcommand{\disciplina}{Desigualdade, Diversidade e Redes}
\newcommand{\titulo}{Da produção à demanda}
\newcommand{\subtitulo}{Estrutura da rede intersetorial de atribuição de CO\textsubscript{2} no Brasil e implicações para a coordenação da mitigação}
\newcommand{\autor}{Gabriel Pincelli Cintra}
\newcommand{\local}{Florianópolis}
\newcommand{\dataentrega}{18 de outubro de 2026}
\newcommand{\imprimirtitulodois}{
\begin{center}
  {\Large\bfseries\MakeUppercase{\instituicao}\par}
  \vspace{0.2em}{\large \departamento\par}
  \vspace{0.2em}{\normalsize\textbf{\disciplina}\ (\curso)\par}
  \vspace{1.4em}{\LARGE\bfseries \titulo\par}
  \vspace{0.25em}{\large\itshape \subtitulo\par}
  \vspace{1.5em}
\end{center}
\noindent\begin{minipage}{0.48\textwidth}
  \raggedright\normalsize \autor
\end{minipage}\hfill
\begin{minipage}{0.48\textwidth}
  \raggedleft\normalsize \local, \dataentrega
\end{minipage}
\vspace{1.0em}\hrule\vspace{1.0em}}
\begin{document}
\imprimirtitulodois
% Resumo e palavras-chave serão redigidos após os resultados.
\input{introducao.tex}
% Sequência inspirada em Huang et al. (2019).
\section{Dados e recorte da análise}
\label{sec:dados}

\subsection{Fontes, unidades e preparação}

A base econômica é a MIP brasileira de 2015, disponibilizada pelo IBGE no detalhamento de 67 atividades \parencite{ibge2015mip}. A unidade da análise de redes é a atividade produtiva. O recorte nacional de um único ano não distingue unidades da Federação nem permite examinar mudanças temporais. Os valores monetários são expressos em milhões de reais, a preços básicos de 2015.

A Tabela~\ref{tab:fontes} reúne os insumos econômicos e ambientais. As intensidades são expressas em gigagramas de CO\textsubscript{2} por milhão de reais (Gg CO\textsubscript{2}/R\$ milhão). A demanda final por produtos nacionais agrega consumo das famílias, governo e instituições sem fins de lucro, formação bruta de capital fixo, variação de estoques e exportações. O vetor por atividade é obtido distribuindo a demanda dos 127 produtos entre suas atividades produtoras, conforme as participações da Tabela 13 do IBGE. As categorias de uso final não serão representadas por nós separados.

\begin{table}[htbp]
\centering
\caption{Fontes e objetos de entrada da análise.}
\label{tab:fontes}
\small
\begin{tabular}{@{}>{\raggedright\arraybackslash}p{0.26\textwidth}>{\raggedright\arraybackslash}p{0.19\textwidth}>{\raggedright\arraybackslash}p{0.46\textwidth}@{}}
\toprule
Objeto & Dimensão & Fonte e unidade \\
\midrule
Produção bruta ($x$) & $67\times1$ & IBGE, Tabela 01; R\$ milhões \\
Demanda final por atividade ($y$) & $67\times1$ & IBGE, Tabela 03, convertida com a Tabela 13; R\$ milhões \\
Coeficientes técnicos ($A$) & $67\times67$ & IBGE, Tabela 14; adimensional \\
Inversa de Leontief ($L$) & $67\times67$ & IBGE, Tabela 15; adimensional \\
Intensidades diretas ($\varphi^{2011}$, $\varphi^{2018}$) & $67\times1$ por ano & Sanguinet e Azzoni (2024), Tabela 2; Gg CO\textsubscript{2}/R\$ milhão \\
\bottomrule
\end{tabular}
\par\smallskip
\begin{minipage}{\textwidth}\footnotesize
Fonte: IBGE (2018) e Sanguinet e Azzoni (2024). Os anos ambientais são 2011 e 2018; a base econômica refere-se a 2015. As dimensões descrevem a organização das entradas, não observações estatisticamente independentes.
\end{minipage}
\end{table}

A Tabela 2 de \textcite{sanguinetAzzoni2024} define a intensidade como emissão direta de CO\textsubscript{2} por unidade de produção bruta. Os 134 valores dos dois anos foram transcritos e conferidos com a fonte. A correspondência nominal com as 67 atividades da MIP é direta, sem agregação ou desagregação local; isso não valida a alocação ambiental anterior realizada pelos autores. A tabela fornece um vetor por ano, sem intensidades observadas separadamente para cada estado.

Para o cenário de 2015, utiliza-se interpolação linear:
\begin{equation}
\varphi_i^{\mathrm{cen},2015}=\frac{3}{7}\varphi_i^{2011}+\frac{4}{7}\varphi_i^{2018}.
\label{eq:coeficiente-cenario}
\end{equation}
Os pesos decorrem da posição de 2015 no intervalo. A interpolação pressupõe evolução linear da intensidade e não incorpora observação física de emissões em 2015. Como a fonte publica duas casas decimais, casas adicionais do cálculo não representam precisão ambiental adicional.

\subsection{Fronteira e validade ambiental}

A fronteira pretendida abrange emissões de CO\textsubscript{2} das atividades produtivas brasileiras representadas na MIP, para demanda doméstica e exportações. Exclui emissões estrangeiras na produção de importações e nas etapas posteriores de processamento ou uso dos produtos exportados. O destino da atribuição identifica o setor do produto final na fronteira nacional, não seu país comprador ou a residência do consumidor. A conta pela demanda não constitui uma pegada mundial do consumo dos residentes.

Os resultados são expressos em Gg de CO\textsubscript{2}, com eventual conversão por $1\,000$ Gg $=1$ Mt. Não se contabiliza o conjunto dos gases de efeito estufa em CO\textsubscript{2} equivalente. A cobertura de combustão, processos industriais, biomassa e mudanças de uso da terra depende das categorias da fonte. A auditoria documental encontrou descrições divergentes sobre uso da terra e não recuperou os numeradores físicos nem os procedimentos completos de alocação necessários para reconstruir as intensidades.

A compatibilidade monetária também permanece pendente. A fonte relata atualização de sua matriz inter-regional para preços de dezembro de 2018, enquanto a aplicação exploratória utiliza os coeficientes interpolados com valores de 2015, sem ajuste adicional. A igualdade das unidades não resolve a diferença de preços; um fator comum não substitui a verificação dos denominadores setoriais. Comparações preliminares com informações de processos industriais apontaram discrepâncias de magnitude, documentadas na auditoria do repositório\footnote{Documentação: \texttt{docs/auditoria\_coeficientes\_emissoes\_2015.md}. As entradas e seus identificadores SHA-256 constam de \texttt{raw/manifesto.csv}.}.

A base ambiental atual é, portanto, um \textbf{cenário exploratório}. Sua utilização em conclusões empíricas definitivas depende de reconstrução e validação ou de substituição por uma extensão compatível com 2015. As verificações de balanço e das identidades contábeis asseguram consistência dos cálculos, mas não certificam a validade das emissões estimadas.\footnote{A documentação computacional reúne notebooks exploratórios e versões de trabalho. As definições desta análise são as apresentadas neste texto; resultados dos notebooks não são incorporados automaticamente ao trabalho.}

\section{Metodologia}
\subsection{Da MIP à matriz de atribuição de emissões}
\label{sec:derivacao}

A análise parte da matriz $A$ de coeficientes técnicos de insumos produzidos no Brasil, da produção bruta $x$ e da demanda final por atividade $y$ da MIP do IBGE. Os índices $i$ e $j$ identificam atividades e os vetores são colunas. Cada $A_{ij}$ indica o valor do insumo produzido em $i$ requerido por unidade monetária de produção de $j$. O balanço e a inversa de Leontief são
\begin{equation}
x=Ax+y=Ly,\qquad L=(I-A)^{-1}.
\label{eq:leontief}
\end{equation}
$L_{ij}$ reúne a produção de $i$ requerida direta e indiretamente por unidade de demanda final de $j$, sob coeficientes técnicos fixos \parencite{millerBlair2009}. Utilizam-se $A$ e $L$ das Tabelas 14 e 15, com conferência de consistência e do balanço $Ly=x$ considerando o arredondamento da fonte.

Para incorporar a dimensão ambiental, seja $\varphi_i$ a intensidade direta de CO\textsubscript{2} por unidade de produção de $i$, em Gg de CO\textsubscript{2} por milhão de reais, com a origem e as limitações descritas na seção~\ref{sec:dados}. Organizando essas intensidades em $\Phi=\operatorname{diag}(\varphi)$, obtém-se
\begin{equation}
C=\Phi L,\qquad C_{ij}=\varphi_iL_{ij}.
\label{eq:coeficientes-emissoes}
\end{equation}
Cada $C_{ij}$ expressa a emissão em $i$ por unidade monetária de demanda final de $j$ e tem unidade de Gg de CO\textsubscript{2} por milhão de reais. Essa extensão ambiental relaciona as intensidades diretas aos requisitos da cadeia, conforme a abordagem utilizada por \textcite{sanguinetAzzoni2024}.

Multiplicar cada coluna de $C$ pela demanda final efetiva fornece a matriz de atribuição:
\begin{equation}
F=C\widehat y=\Phi L\widehat y,\qquad F_{ij}=\varphi_iL_{ij}y_j,
\label{eq:matriz-atribuicao}
\end{equation}
em que $\widehat y=\operatorname{diag}(y)$. $F$ tem dimensão $67\times67$ e unidade de Gg de CO\textsubscript{2}. O elemento $F_{ij}$ representa a emissão ocorrida em $i$ e atribuída à demanda final dos produtos de $j$, incluindo requisitos diretos e indiretos. Essa é a matriz que conecta as duas atribuições discutidas por \textcite{marques2012} e fundamenta a representação em rede adotada a partir de \textcite{huang2019}. Seus elementos representam atribuição contábil, não transações diretas nem transferências físicas de CO\textsubscript{2}.

\subsection{Atribuição pela produção e pela demanda}
\label{sec:atribuicoes}

A matriz $F$ permite organizar as mesmas emissões segundo a atividade em que ocorrem ou o produto final ao qual são atribuídas. A Figura~\ref{fig:exemplo-atribuicoes} ilustra essa distinção em uma economia de quatro atividades: Eletricidade, Agricultura, Alimentos e Restaurantes.

Na atribuição pela produção, somam-se as linhas de $F$, preservando a origem emissora:
\begin{equation}
c_i=\sum_jF_{ij}=\varphi_i x_i,\qquad c=F\mathbf{1}=\Phi x.
\label{eq:conta-producao}
\end{equation}
Aqui, $\mathbf{1}$ é um vetor-coluna de uns. No quadro (a), as emissões da geração elétrica são contabilizadas em Eletricidade, inclusive quando sua produção atende às outras atividades.

Na atribuição pela demanda, somam-se as colunas de $F$, preservando o setor do produto final:
\begin{equation}
u_j=\sum_iF_{ij}=y_j\sum_i\varphi_iL_{ij},\qquad u=F^\top\mathbf{1},
\label{eq:conta-demanda}
\end{equation}
em que $\top$ indica transposição. No quadro (b), a atribuição à demanda final de Restaurantes reúne as emissões de sua própria atividade e as de Eletricidade, Agricultura e Alimentos associadas ao atendimento dessa demanda. As transações e a produção são as mesmas do quadro (a); muda o critério de atribuição representado pelas cores.

O Sankey do quadro (c) conecta essas perspectivas: cada faixa representa $F_{ij}$, a barra de origem totaliza $c_i$ e a de destino totaliza $u_j$. Por exemplo, a faixa entre Eletricidade e Restaurantes corresponde à emissão da primeira atribuída à demanda final da segunda. Enquanto as setas superiores representam transações monetárias, as faixas do Sankey representam atribuições de emissões. O total é preservado:
\begin{equation}
\mathbf{1}^\top c=\mathbf{1}^\top F\mathbf{1}=\mathbf{1}^\top u.
\label{eq:conservacao-emissoes}
\end{equation}

Para uma mesma atividade, a diferença entre as contas é o saldo das atribuições entre atividades distintas:
\begin{equation}
u_i-c_i=\sum_{j\ne i}F_{ji}-\sum_{j\ne i}F_{ij}.
\label{eq:saldo-atribuicao}
\end{equation}
A diagonal cancela nesse saldo, mas permanece nas contas completas. A diferença pode resultar de poucas relações de grande peso ou de muitas relações distribuídas, o que motiva examinar sua concentração e conectividade na rede.

% Figura vetorial copiada sem edição de:
% C:/GitHub/Mestrado/0-Projeto de Dissertação/figura_economia/contabilidade_com_sankey.pdf
% SHA-256: 861F8B5C1FDEBE3C96B31D8BD96763C2829FA2AE72703160C83949E4ACD6BF29
\begin{figure}[p]
\centering
\includegraphics[width=\textwidth,height=0.76\textheight,keepaspectratio]{figuras/atribuicao_cadeia_alimentos.pdf}
\caption{Atribuição das emissões pela produção e pela demanda em uma economia ilustrativa.}
\label{fig:exemplo-atribuicoes}
\begin{minipage}{\textwidth}\footnotesize
Fonte: elaboração própria.

Nota: exemplo didático, sem representar estimativas nacionais dos quatro setores. Em (a--b), áreas e espessuras das setas expressam valores monetários; no nó Demanda Final, a área indica o valor recebido, incluindo exportações. As cores representam $c$ em (a) e $u$ em (b). Em (c), as barras à esquerda representam as origens emissoras (I--IV) e as barras à direita, os produtos finais (V--VIII). Suas alturas e as espessuras das faixas expressam emissões em Gg de CO\textsubscript{2}. Os três quadros compartilham a escala de cores.
\end{minipage}
\end{figure}
\clearpage

\subsection{Construção da rede dirigida e ponderada}
\label{sec:construcao-rede}

A rede terá 67 nós, um por atividade da MIP, identificados pelos códigos do IBGE. Cada nó representa uma origem emissora nas ligações de saída e o setor de um produto final nas ligações de entrada, seguindo a representação de atribuição de \textcite{huang2019}. Os dois papéis pertencem ao mesmo nó, enquanto o Sankey da Figura~\ref{fig:exemplo-atribuicoes} os apresenta em conjuntos separados.

Para estudar relações entre atividades distintas, define-se
\begin{equation}
W_{ij}=\begin{cases}
F_{ij}, & i\ne j,\\
0, & i=j.
\end{cases}
\label{eq:rede-sem-diagonal}
\end{equation}
Haverá uma aresta $i\rightarrow j$ para cada $W_{ij}>0$, com peso igual a esse valor, em Gg de CO\textsubscript{2}. A direção indica a origem emissora e o destino da atribuição, e as duas direções de um par serão tratadas separadamente. Como os pesos já incorporam requisitos diretos e indiretos, concatenar essas arestas não representa uma sequência de etapas produtivas.

A diagonal $F_{ii}$ reúne emissões de $i$ atribuídas à demanda final de seus próprios produtos, incluindo os requisitos indiretos correspondentes. Ela permanece nas contas $c$ e $u$, mas é retirada de $W$ para separar as relações entre atividades da atribuição à própria atividade. Assim, $W$ não terá laços, e seu peso total corresponderá apenas à parcela intersetorial das emissões.

Serão mantidos todos os nós e todas as arestas positivas, sem corte por magnitude nesta etapa. A construção exige alinhamento dos códigos e pesos finitos e não negativos; eventuais valores negativos deverão ser examinados antes da representação, sem conversão automática para zero ou valor absoluto. A seleção de arestas será tratada na seção~\ref{sec:backbone}.

\subsection{Volume e concentração das relações}
\label{sec:volume-concentracao}

Para relacionar a estrutura da rede às duas contas ambientais, serão examinados o volume das relações externas e sua distribuição entre contrapartes. As medidas serão calculadas em $W$, antes da filtragem, com referência às forças de \textcite{huang2019} e à disparidade dos pesos de \textcite{serrano2009}.

\paragraph{Volume.}
As forças de saída e entrada são as somas dos pesos em cada direção:
\begin{equation}
\begin{aligned}
s_i^{\mathrm{out}}&=\sum_{j\ne i}W_{ij}=c_i-F_{ii},\\
s_i^{\mathrm{in}}&=\sum_{j\ne i}W_{ji}=u_i-F_{ii}.
\end{aligned}
\label{eq:forcas-externas}
\end{equation}
A saída mede as emissões próprias de $i$ atribuídas aos produtos finais de outras atividades; a entrada mede as emissões de outras atividades atribuídas aos produtos finais de $i$. Ambas têm unidade de Gg de CO\textsubscript{2}. Sua diferença é $u_i-c_i$, mas a leitura separada permite identificar grandes volumes que se compensam no saldo. Serão apresentados também $c_i$, $u_i$ e $F_{ii}$, para situar a parcela intersetorial em cada conta completa.

\paragraph{Concentração.}
Para perfis com força positiva, define-se a participação de cada contraparte:
\begin{equation}
p_{ij}^{\mathrm{out}}=\frac{W_{ij}}{s_i^{\mathrm{out}}},\qquad
p_{ji}^{\mathrm{in}}=\frac{W_{ji}}{s_i^{\mathrm{in}}},\quad j\ne i.
\label{eq:participacoes-externas}
\end{equation}
O perfil de saída distribui as emissões externas da origem entre destinos; o de entrada descreve as origens que compõem a atribuição externa ao destino. Cada perfil soma um. A disparidade e seu inverso serão calculados por
\begin{equation}
\begin{aligned}
Y_i^{\mathrm{out}}&=\sum_{j\ne i}(p_{ij}^{\mathrm{out}})^2,&
Y_i^{\mathrm{in}}&=\sum_{j\ne i}(p_{ji}^{\mathrm{in}})^2,\\
D_{2,i}^{\mathrm{out}}&=1/Y_i^{\mathrm{out}},&
D_{2,i}^{\mathrm{in}}&=1/Y_i^{\mathrm{in}}.
\end{aligned}
\label{eq:disparidade}
\end{equation}
$Y$ aumenta com a concentração; $D_2$ expressa a mesma informação como número efetivo de contrapartes. Para $k$ pesos iguais, $D_2=k$; quando uma relação domina, aproxima-se de um. Essas medidas são adimensionais e ficam indefinidas quando a força é zero, sem excluir o nó da análise.

A combinação de força e número efetivo permitirá distinguir volumes elevados concentrados em poucas contrapartes de volumes distribuídos entre muitas. As atividades que compõem cada perfil serão identificadas pelos pesos e participações, pois o número efetivo não informa sua identidade. O grau contará as arestas positivas de entrada ou saída; a densidade será o número de arestas dividido por $67\times66$. Grau e densidade descrevem a presença de relações, enquanto a disparidade considera a distribuição de seus pesos.

\subsection{Extração e avaliação do backbone}
\label{sec:backbone}

O filtro de disparidade de \textcite{serrano2009}, aplicado por \textcite{huang2019}, será utilizado para selecionar ligações destacadas na composição local dos pesos. A seleção depende da participação da aresta na força de cada extremidade, permitindo considerar atividades de diferentes escalas. Seu resultado será o backbone, sem alterar as contas ambientais completas.

\paragraph{Critério de seleção.}
O modelo nulo divide o intervalo unitário em $k$ parcelas por $k-1$ pontos escolhidos uniformemente, representando uma distribuição aleatória da força entre as ligações. Para $k>1$, a probabilidade de uma parcela atingir ou superar $p$ é $(1-p)^{k-1}$. Para a aresta $i\rightarrow j$, as caudas nas perspectivas de saída e entrada são
\begin{equation}
\begin{aligned}
a_{ij}^{\mathrm{out}}&=\left(1-\frac{W_{ij}}{s_i^{\mathrm{out}}}\right)^{k_i^{\mathrm{out}}-1},\\
a_{ij}^{\mathrm{in}}&=\left(1-\frac{W_{ij}}{s_j^{\mathrm{in}}}\right)^{k_j^{\mathrm{in}}-1}.
\end{aligned}
\label{eq:caudas-filtro}
\end{equation}
Os graus $k$ e as forças serão os de $W$. A aresta será selecionada se $a_{ij}^{\mathrm{out}}<\alpha$ \textbf{ou} $a_{ij}^{\mathrm{in}}<\alpha$, registrando-se qual extremidade determina a seleção. A saída destaca um destino nas atribuições da origem; a entrada destaca uma origem na composição do destino.

As caudas só serão calculadas para graus pertinentes maiores que um. Se apenas uma extremidade tiver grau um, valerá o teste da outra; se ambas tiverem grau um, a aresta será preservada por convenção de manutenção da relação exclusiva, inclusive em díades isoladas. Essa extensão dos casos degenerados do suplemento de Serrano et al. será registrada separadamente da seleção estatística. Todos os nós serão mantidos.

Será adotado $\alpha=0{,}05$, com $0{,}01$ e $0{,}10$ para sensibilidade. Para cada valor, o filtro será aplicado uma vez a $W$, sem atualização iterativa dos perfis. A matriz resultante $B_\alpha$ manterá os pesos originais nas arestas selecionadas e zero nas demais, sem renormalização ou corte absoluto adicional.

\paragraph{Avaliação da seleção.}
Serão comparadas a fração de arestas retidas e a cobertura dos pesos, globalmente e por atividade:
\begin{equation}
\begin{aligned}
r_\alpha&=\frac{\sum_{ij}(B_\alpha)_{ij}}{\sum_{ij}W_{ij}},\\
r_{\alpha,i}^{\mathrm{out}}&=\frac{\sum_j(B_\alpha)_{ij}}{s_i^{\mathrm{out}}},\qquad
r_{\alpha,i}^{\mathrm{in}}=\frac{\sum_j(B_\alpha)_{ji}}{s_i^{\mathrm{in}}}.
\end{aligned}
\label{eq:cobertura-backbone}
\end{equation}
As razões serão definidas apenas para denominadores positivos. A cobertura se refere à parcela intersetorial de $W$; sua leitura por atividade permite identificar perdas ocultas pela retenção agregada. Serão descritos também graus, forças retidas, isolados e componentes fracamente conectados, obtidos ao desconsiderar a direção das arestas para verificar conectividade.

A sensibilidade será examinada pelas mudanças nas arestas, coberturas e posições setoriais entre os três valores de $\alpha$. Com convenções fixas, os conjuntos são aninhados por construção; essa propriedade não constitui evidência independente de robustez. O modelo nulo é uma referência de distribuição dos pesos: suas caudas não medem a incerteza ambiental, e a seleção não estabelece prioridade de mitigação.

A resposta à pergunta combinará as contas e seus saldos em $F$, os volumes e a concentração em $W$ e as contrapartes destacadas em $B_\alpha$. As relações serão interpretadas junto de seus pesos e coberturas: uma ligação pode ser importante localmente e pequena no total, enquanto uma ligação não selecionada continua integrando a contabilidade. Assim, o backbone complementa a análise dos perfis completos, sem substituí-los.

\section{Resultados}
\label{sec:resultados}

A análise identifica uma rede intersetorial quase completa, mas com distribuições de pesos e papéis de atribuição distintos entre atividades. A filtragem reduz o número de ligações e evidencia essas diferenças, com perdas de cobertura desiguais. Os valores apresentados correspondem ao cenário ambiental exploratório descrito na seção~\ref{sec:dados}; não constituem um inventário validado das emissões brasileiras de 2015.

\subsection{Contas ambientais, volume e concentração}

O total atribuído é de 683,70 mil Gg de CO\textsubscript{2} nas duas contas. A diagonal de $F$ reúne 392,92 mil Gg, e as relações entre atividades distintas somam 290,78 mil Gg em $W$. A igualdade entre os totais pela produção e pela demanda decorre da contabilidade; o resultado de interesse é como as mesmas emissões se distribuem entre atividades e contrapartes.

A rede completa tem 67 nós e 4\,290 arestas, com densidade de 97,01\%. As 66 atividades com relações externas estão ligadas entre si nas duas direções e têm graus de entrada e saída iguais a 65. Serviços domésticos (9700) é o único isolado: sua conta permanece na diagonal de $F$. Assim, os graus da rede completa pouco distinguem os perfis, enquanto os pesos mostram diferenças de volume e concentração.

A Figura~\ref{fig:indicadores-resultados} compara essas dimensões para as atividades que compõem a união das seis maiores forças de saída e das seis maiores forças de entrada. Energia (3500) e água, esgoto e resíduos (3680), por exemplo, distribuem volumes externos próximos: 31,32 e 31,23 mil Gg. Entretanto, seus números efetivos de destinos são 20,79 e 9,88. No cenário analisado, volumes semelhantes correspondem a atribuições mais distribuídas em energia e mais concentradas em água e resíduos. Entre os 66 perfis definidos, a mediana é de 11,50 destinos efetivos e 13,93 origens efetivas, abaixo das 65 contrapartes existentes em cada direção.

A Tabela~\ref{tab:contas-energia-construcao} situa energia e construção nas contas completas. Em energia, a produção supera a atribuição pela demanda em 28,39 mil Gg; em construção ocorre o inverso, por 20,56 mil Gg. Os volumes externos mostram como esses saldos se compõem: energia distribui 31,32 mil Gg para outros produtos finais e recebe 2,92 mil Gg de outras origens; construção distribui 2,93 mil Gg e recebe 23,48 mil Gg. A pequena saída externa de construção não indica baixa emissão própria: 35,91 dos seus 38,84 mil Gg pela produção são atribuídos à demanda final da própria atividade e ficam fora de $W$.

\begin{table}[htbp]
\centering
\caption{Contas e volumes intersetoriais de energia e construção.}
\label{tab:contas-energia-construcao}
\small
\begin{tabular}{@{}lrr@{}}
\toprule
Medida & Energia (3500) & Construção (4180) \\
\midrule
Produção, $c$ & 56,62 & 38,84 \\
Demanda, $u$ & 28,23 & 59,40 \\
Atribuição à própria atividade, $F_{ii}$ & 25,30 & 35,91 \\
Saída externa, $s^{\mathrm{out}}$ & 31,32 & 2,93 \\
Entrada externa, $s^{\mathrm{in}}$ & 2,92 & 23,48 \\
Saldo, $u-c$ & $-28,39$ & 20,56 \\
\bottomrule
\end{tabular}
\par\smallskip\begin{minipage}{\textwidth}\footnotesize
Fonte: elaboração própria. Valores em mil Gg de CO\textsubscript{2}, antes do filtro. A diagonal integra $c$ e $u$, mas não as forças externas. Diferenças residuais entre somas de valores exibidos decorrem do arredondamento.
\end{minipage}
\end{table}

\clearpage
\newgeometry{margin=1.5cm}
\begin{landscape}
\begin{figure}[p]
\centering
\includegraphics[width=\linewidth,height=0.80\textheight,keepaspectratio]{figuras/indicadores.pdf}
\caption{Volume, concentração, conectividade selecionada e cobertura em atividades de maior volume externo.}
\label{fig:indicadores-resultados}
\begin{minipage}{\linewidth}\footnotesize
Fonte: elaboração própria. Linhas: união das seis maiores forças de saída e das seis maiores forças de entrada de $W$; não representam todas as atividades. Azul: saída; laranja: entrada. Volume e número efetivo são calculados em $W$ completa. Grau e cobertura referem-se a $\alpha=0{,}05$; os intervalos do grau abrangem $\alpha=0{,}01$--$0{,}10$ e não são intervalos de confiança. A diagonal é excluída.
\end{minipage}
\end{figure}
\end{landscape}
\restoregeometry

\subsection{Seleção das relações e limites de cobertura}

O backbone de referência preserva 575 arestas, ou 13,40\% das ligações de $W$, e 70,69\% de seu peso intersetorial (Tabela~\ref{tab:sensibilidade-backbone}). A matriz de adjacência da Figura~\ref{fig:matriz-backbone} mostra essas mesmas relações pela presença e pelo peso. A linha extensa de energia indica muitos destinos retidos; a coluna extensa de construção indica muitas origens retidas. A comparação entre os painéis distingue a existência de uma ligação de seu volume: nem todas as células retidas carregam pesos elevados.

\begin{table}[htbp]
\centering
\caption{Sensibilidade da seleção e da cobertura ao parâmetro do filtro.}
\label{tab:sensibilidade-backbone}
\small
\begin{tabular}{@{}crrrr@{}}
\toprule
$\alpha$ & Arestas & \% das arestas de $W$ & Peso retido & Cobertura (\%) \\
\midrule
0,01 & 338 & 7,88 & 166,88 & 57,39 \\
0,05 & 575 & 13,40 & 205,55 & 70,69 \\
0,10 & 755 & 17,60 & 224,92 & 77,35 \\
\bottomrule
\end{tabular}
\par\smallskip\begin{minipage}{\textwidth}\footnotesize
Fonte: elaboração própria. Peso em mil Gg de CO\textsubscript{2}. Cobertura relativa ao peso de $W$, sem diagonal. Todos os casos mantêm 67 nós, um isolado e um componente fracamente conectado de 66 atividades. O filtro é aplicado à mesma rede completa, com regra OU e sem atualização dos perfis.
\end{minipage}
\end{table}

A cobertura agregada oculta diferenças locais. Em $\alpha=0{,}05$, as medianas entre os perfis definidos são 59,32\% nas saídas e 55,69\% nas entradas; os mínimos são 32,96\% e 33,79\%. Portanto, a preservação de parcela expressiva do peso total não assegura representação igualmente completa de cada atividade e direção.

A sensibilidade também limita a interpretação dos graus. Em relação a $\alpha=0{,}05$, a redução para 0,01 altera os postos de grau de saída em mediana de 11 posições, e os de entrada em sete; a elevação para 0,10 produz mudanças medianas de 4,5 e quatro posições. Os postos consideram empates. A permanência do componente de 66 atividades nos três casos descreve a conectividade agregada, mas não demonstra estabilidade das posições individuais. Mantém-se 0,05 como referência, sem tratá-lo como parâmetro ótimo.

Dos 575 vínculos retidos, 232 passam somente pelo teste de saída na origem, 265 somente pelo de entrada no destino e 78 pelos dois. Nenhum depende da convenção de grau um. Essa composição mostra por que o backbone não equivale a selecionar apenas as maiores emissões absolutas: a relação pode destacar-se em uma ponta e representar parcela pequena da força da outra.

\clearpage
\newgeometry{margin=1.5cm}
\begin{landscape}
\begin{figure}[p]
\centering
\includegraphics[width=\linewidth,height=0.82\textheight,keepaspectratio]{figuras/matriz_presenca.pdf}
\caption{Matriz de adjacência do backbone com $\alpha=0{,}05$: (a) presença das arestas.}
\label{fig:matriz-backbone}
\begin{minipage}{\linewidth}\footnotesize
Fonte: elaboração própria. Linha: origem emissora; coluna: atividade do produto final. Ordem setorial do IBGE nos dois eixos. Azul identifica as 575 ligações retidas; cinza indica ausência no backbone, inclusive a diagonal. Ausência no backbone não implica valor nulo em $F$ ou $W$. Todas as 67 atividades permanecem nos eixos.
\end{minipage}
\end{figure}
\clearpage
\begin{figure}[p]
\ContinuedFloat
\centering
\includegraphics[width=\linewidth,height=0.82\textheight,keepaspectratio]{figuras/matriz_peso.pdf}
\caption{Matriz de adjacência do backbone com $\alpha=0{,}05$: (b) pesos das mesmas arestas.}
\begin{minipage}{\linewidth}\footnotesize
Fonte: elaboração própria. Mesma ordem e mesma seleção do painel (a). A cor representa pesos originais em Gg de CO\textsubscript{2}, em escala logarítmica, sem normalização por linha ou coluna. Cinza identifica células sem aresta retida. A cor mede volume absoluto; a concentração é descrita pelos números efetivos em $W$ completa.
\end{minipage}
\end{figure}
\end{landscape}
\restoregeometry

\subsection{Energia e construção: perspectivas de uma mesma relação}

Energia e construção ilustram dois papéis na atribuição intersetorial: uma origem emissora associada a muitos destinos finais e um destino que reúne emissões de muitas origens. A Figura~\ref{fig:rede-focos} apresenta duas vistas completas do mesmo backbone, com posições idênticas e destaque para as relações de cada atividade. Os 33 nós com $c>u$ ficam à esquerda, os 33 com $u>c$ à direita, e o isolado fica abaixo. Essa organização explicita o saldo já calculado; não resulta de detecção de comunidades. As cores mostram entradas e saídas do destaque, enquanto os motivos de retenção são detalhados na Tabela~\ref{tab:motivos-casos} e nos painéis completos do Anexo~\ref{anexo:paineis}.

\begin{table}[htbp]
\centering
\caption{Motivos de retenção e cobertura nos recortes de energia e construção.}
\label{tab:motivos-casos}
\small
\begin{tabular}{@{}lrrrrr@{}}
\toprule
Recorte & Arestas & Só saída & Só entrada & Ambos & Cobertura (\%) \\
\midrule
Energia: entradas & 7 & 1 & 5 & 1 & 42,1 \\
Energia: saídas & 62 & 0 & 56 & 6 & 99,2 \\
Construção: entradas & 45 & 41 & 0 & 4 & 91,4 \\
Construção: saídas & 5 & 4 & 1 & 0 & 57,0 \\
\bottomrule
\end{tabular}
\par\smallskip\begin{minipage}{\textwidth}\footnotesize
Fonte: elaboração própria. $\alpha=0{,}05$. Saída e entrada nas colunas de motivo identificam o teste na origem e no destino da aresta, independentemente da atividade em destaque. Cobertura: peso retido dividido pela força da mesma atividade e direção em $W$. Os recortes compartilham arestas e não devem ser somados para obter totais da rede.
\end{minipage}
\end{table}

Em energia, 56 das 62 saídas retidas passam somente pelo teste de entrada no destino. A atividade tem participação saliente na composição externa de muitas demandas finais, embora cada destino possa representar parcela pequena das emissões distribuídas pela origem. As outras seis relações passam pelos dois testes. Entre os maiores destinos estão comércio, administração pública, construção e atividades alimentares.

Em construção, 41 das 45 entradas passam somente pelo teste de saída na origem. Construção representa parcela saliente das atribuições externas dessas origens, mesmo quando a contribuição individual é pequena no total recebido pelo destino. As quatro entradas que passam pelos dois testes vêm de produtos de minerais não metálicos, energia, água e resíduos e comércio. A relação energia $\rightarrow$ construção, comum aos dois recortes, pesa aproximadamente 2\,263 Gg: 7,2\% das saídas externas de energia e 9,6\% das entradas externas de construção.

As direções predominantes também apresentam maior cobertura: 99,2\% das saídas de energia e 91,4\% das entradas de construção. Nas direções inversas, ficam fora do backbone 57,9\% das entradas externas de energia e 43,0\% das saídas externas de construção. Assim, os recortes representam bem os volumes predominantes, mas as listas curtas nas outras direções não autorizam desprezar os vínculos omitidos em conjunto. A diferença entre 62 saídas de energia e 45 entradas de construção também não mede concentração: em $W$ completa, os respectivos números efetivos são próximos, ambos de 20,8 contrapartes.

Em conjunto, os resultados mostram relações de pesos desiguais cuja saliência depende da ponta examinada. Forças e números efetivos caracterizam os perfis completos; o backbone identifica contrapartes com cobertura conhecida. O saldo $u-c$ permanece calculado nas contas completas. As arestas descrevem atribuições diretas e indiretas; sua concatenação não representa trajetórias produtivas.

\clearpage
\begin{figure}[p]
\centering
\includegraphics[width=\textwidth,height=0.76\textheight,keepaspectratio]{figuras/rede_foco_3500.pdf}
\caption{Duas perspectivas do mesmo backbone: (a) destaque para energia (3500).}
\label{fig:rede-focos}
\begin{minipage}{\textwidth}\footnotesize
Fonte: elaboração própria. Cada vista mantém 67 nós e 575 arestas em $\alpha=0{,}05$. Azul: saídas do destaque; laranja: entradas; cinza: demais arestas. Energia tem 62 saídas e sete entradas, com cobertura de 99,2\% e 42,1\% do volume externo de cada direção. Área dos nós uniforme; espessura cresce com a raiz quadrada do peso. Posição definida pelo sinal de $u-c$ nas contas completas e pela ordem IBGE em cada lado, não por centralidade ou comunidade.
\end{minipage}
\end{figure}
\clearpage
\begin{figure}[p]
\ContinuedFloat
\centering
\includegraphics[width=\textwidth,height=0.76\textheight,keepaspectratio]{figuras/rede_foco_4180.pdf}
\caption{Duas perspectivas do mesmo backbone: (b) destaque para construção (4180).}
\begin{minipage}{\textwidth}\footnotesize
Fonte: elaboração própria. Mesmos nós, arestas, posições e escala de espessura do painel (a). Construção tem cinco saídas e 45 entradas, com cobertura de 57,0\% e 91,4\%. Cores representam direção em relação ao destaque, não motivo de retenção. As duas vistas preservam 70,69\% do peso intersetorial de $W$. Cruzamentos e proximidade não representam trajetórias ou distâncias econômicas; relações individuais podem ser consultadas na matriz e nos painéis do anexo.
\end{minipage}
\end{figure}
\clearpage

\section{Discussão e implicações para a coordenação da mitigação}
Intervenções nos processos produtivos e no uso de energia e materiais ao longo das cadeias constituem frentes complementares de mitigação \parencite{rissman2020}.

% Reservado: interpretação conjunta e limites dos dados e métodos.
\section{Conclusões}
% Reservado até a conclusão da análise.
\printbibliography[title={Referências}]
\clearpage
\appendix
\section{Painéis completos de energia e construção}
\label{anexo:paineis}

Os painéis apresentam todas as arestas dos recortes em $\alpha=0{,}05$, com pesos originais em Gg de CO\textsubscript{2} e escala linear comum, sem novo filtro. Cobertura: peso retido sobre a força da atividade e direção em $W$, sem diagonal. Azul indica somente o teste de saída na origem; laranja, somente entrada no destino; roxo, ambos. A paginação preserva a seleção e os denominadores.

\begin{figure}[htbp]
\centering
\includegraphics[width=\textwidth,height=0.85\textheight,keepaspectratio]{figuras/anexo_energia_1.pdf}
\caption{Energia (3500): todas as sete entradas e as saídas 1--31 de 62.}
\label{fig:anexo-energia}
\begin{minipage}{\textwidth}\footnotesize
Fonte: elaboração própria. Coberturas de 42,1\% nas entradas e 99,2\% no conjunto das saídas. As demais saídas são apresentadas na continuação; não há novo filtro. Pesos exibidos com uma casa decimal. Azul: só saída na origem; laranja: só entrada no destino; roxo: ambos os testes.
\end{minipage}
\end{figure}
\clearpage
\begin{figure}[p]
\ContinuedFloat
\centering
\includegraphics[width=\textwidth,height=0.85\textheight,keepaspectratio]{figuras/anexo_energia_2.pdf}
\caption{Energia (3500), continuação: saídas 32--62 de 62.}
\begin{minipage}{\textwidth}\footnotesize
Fonte: elaboração própria. Mesma seleção, ordem, escala e cores do painel anterior. A cobertura de 99,2\% se refere às 62 saídas, não apenas a esta página. Azul: só saída na origem; laranja: só entrada no destino; roxo: ambos os testes.
\end{minipage}
\end{figure}
\clearpage
\begin{figure}[p]
\centering
\includegraphics[width=\textwidth,height=0.85\textheight,keepaspectratio]{figuras/anexo_construcao.pdf}
\caption{Construção (4180): todas as 45 entradas e cinco saídas.}
\label{fig:anexo-construcao}
\begin{minipage}{\textwidth}\footnotesize
Fonte: elaboração própria. Coberturas de 91,4\% nas entradas e 57,0\% nas saídas. Mesma escala de pesos dos painéis de energia. Azul: só saída na origem; laranja: só entrada no destino; roxo: ambos os testes. Arestas compartilhadas não são duplicadas na rede.
\end{minipage}
\end{figure}
\clearpage

\end{document}
